\documentclass[12pt,a4paper]{article}

% --- PAKETE FÜR DEUTSCHE SPRACHE & SCHRIFTART ---
\usepackage[T1]{fontenc}
\usepackage[ngerman]{babel}
\usepackage{lmodern}

% --- CSQUOTES (Vom Biblatex-Paket empfohlen) ---
\usepackage{csquotes}

% --- MATHE, ENUMERATE, FLOAT & LAYOUT ---
\usepackage{amsmath,amssymb}
\usepackage{geometry}
\geometry{a4paper, margin=2.5cm}
\usepackage{parskip}
\usepackage{enumitem}
\usepackage{float}

% --- TCOLORBOX FÜR HIGHLIGHT- & AUFGABENBOXEN ---
\usepackage[most]{tcolorbox}

\newtcolorbox{taskbox}[1][]{
    colback=blue!5!white,
    colframe=blue!60!black,
    fonttitle=\bfseries,
    title={#1},
    arc=2mm,
    boxrule=0.8pt,
    left=4mm, right=4mm, top=3mm, bottom=3mm
}

\newtcolorbox{emailbox}[1][]{
    colback=red!4!white,
    colframe=red!65!black,
    fonttitle=\bfseries,
    title={#1},
    arc=1.5mm,
    boxrule=0.8pt,
    left=4mm, right=4mm, top=3mm, bottom=3mm
}

% --- FARBEN & VERLINKUNGEN ---
\usepackage{xcolor}
\usepackage[colorlinks=true, linkcolor=blue!60!black, urlcolor=blue!60!black, citecolor=blue!60!black]{hyperref}

% --- CLEVEREF (MUSS NACH HYPERREF GELADEN WERDEN) ---
\usepackage[ngerman]{cleveref}
\crefname{figure}{Abbildung}{Abbildungen}
\Crefname{figure}{Abbildung}{Abbildungen}
\crefname{table}{Tabelle}{Tabellen}
\Crefname{table}{Tabelle}{Tabellen}
\crefname{section}{Abschnitt}{Abschnitte}
\Crefname{section}{Abschnitt}{Abschnitte}

% --- LITERATURVERZEICHNIS MIT BIBLATEX & BIBER ---
\sloppy
\usepackage[backend=biber, style=numeric]{biblatex}
\addbibresource{Skript.bib}

% --- GLOSSAR & ABKÜRZUNGSVERZEICHNIS ---
\usepackage[toc, acronym, nonumberlist]{glossaries}
\makeglossaries

% Custom Glossar-Style für optisch ansprechende Formatierung
\newglossarystyle{mycleanstyle}{%
  \setglossarystyle{long}%
  \renewenvironment{theglossary}%
     {\begin{longtable}{@{}p{5.8cm}@{\hspace{0.6cm}}p{9.6cm}@{}}}%
     {\end{longtable}}%
  \renewcommand*{\glossentry}[2]{%
    \textbf{\glsentryname{##1}} & \glsentrydesc{##1}\\[0.9em]%
  }%
  \renewcommand*{\glsgroupskip}{}%
}
\setglossarystyle{mycleanstyle}

% Inkludieren der separaten Abkürzungs- und Glossardateien
% --- ABKÜRZUNGEN (AKRONYME) ---

\newacronym{mfa}{MFA}{Multi-Faktor-Authentisierung}
\newacronym{nist}{NIST}{National Institute of Standards and Technology}
\newacronym{bsi}{BSI}{Bundesamt für Sicherheit in der Informationstechnik}
\newacronym{smtp}{SMTP}{Simple Mail Transfer Protocol}
\newacronym{url}{URL}{Uniform Resource Locator}
\newacronym{pin}{PIN}{Persönliche Identifikationsnummer}
\newacronym{it}{IT}{Informationstechnik}

% --- GLOSSAREINTRÄGE ---

\newglossaryentry{Multi-Faktor-Authentisierung}{
    name={Multi-Faktor-Authentisierung},
    description={Ein Sicherheitsverfahren, bei dem mehrere unabhängige Faktoren zur Authentifizierung kombiniert werden.}
}

\newglossaryentry{passwort}{
    name={Passwort},
    description={Eine geheime Zeichenkette, die zur Authentisierung dient.}
}

\newglossaryentry{entropie}{
    name={Entropie},
    description={Ein Mass aus der Informationstheorie für die Stärke und Zufälligkeit eines Passworts.}
}

\newglossaryentry{bruteforce}{
    name={Brute-Force-Angriff},
    description={Ein automatisierter Angriff, bei dem Programme alle möglichen Zeichenkombinationen systematisch \mbox{durchprobieren}.}
}

\newglossaryentry{dictionaryattack}{
    name={Dictionary Attack (Wörterbuchangriff)},
    description={Ein gezielter Angriff, bei dem Listen mit \mbox{gebräuchlichen} Wörtern, Namen und bekannten Passwörtern systematisch getestet werden.}
}

\newglossaryentry{datenleck}{
    name={Datenleck},
    description={Ein Vorfall, bei dem vertrauliche Daten unabsichtlich oder absichtlich veröffentlicht oder gestohlen werden.}
}

\newglossaryentry{passwortmanager}{
    name={Passwort-Manager},
    description={Software zur sicheren Verwaltung und Generierung von Passwörtern.}
}

\newglossaryentry{passphrase}{
    name={Passphrase},
    description={Eine längere Zeichenkette, oft bestehend aus mehreren zufälligen Wörtern, die als Passwort verwendet wird.}
}

\newglossaryentry{phishing}{
    name={Phishing},
    description={Eine Form des Social Engineering, bei der Angreifer versuchen, über gefälschte E-Mails, Nachrichten oder Webseiten an vertrauliche Daten zu gelangen.}
}

\newglossaryentry{socialengineering}{
    name={Social Engineering},
    description={Die Manipulation von Personen durch psychologische Tricks, um an vertrauliche Informationen zu gelangen.}
}

\newglossaryentry{spearphishing}{
    name={Spear-Phishing},
    description={Ein gezielter Phishing-Angriff, der auf eine spezifische Person oder bestimmte Mitarbeiter zugeschnitten ist.}
}

\newglossaryentry{quishing}{
    name={Quishing},
    description={Eine Phishing-Methode, bei der manipulierte QR-Codes eingesetzt werden, um Opfer auf gefälschte Webseiten zu leiten.}
}


\begin{document}

% --- TITELBLATT ---
\begin{titlepage}
    \pagenumbering{gobble}
    \centering
    \vspace*{2cm}
    
    {\Huge \textbf{Skript: Passwortsicherheit \& Phishing} \par}
    \vspace{1.5cm}
    {\Large \textbf{Kantonsschule im Lee} \par}
    \vspace{2cm}
    
    \rule{\linewidth}{0.5mm}
    \vspace{2cm}
    
    {\large \textbf{Thema:} Passwortsicherheit, Phishing, Entropie \& Best Practices \par}
    \vspace{0.5cm}
    {\large \textbf{Fach:} Informatik \par}
    \vspace{0.5cm}
    {\large \textbf{Autor:} Laurin Ott \par}
    
    \vfill
    
    {\large \textbf{Datum:} \today \par}
\end{titlepage}

% --- INHALTSVERZEICHNIS ---
\newpage
\pagenumbering{arabic}
\tableofcontents
\newpage

% --- TEIL 1: PASSWORTSICHERHEIT ---
\section{Passwortsicherheit}
\label{sec:passwortsicherheit}

\subsection{Einleitung}
In fast allen IT-Systemen bilden Passwörter die erste Hürde gegen unbefugten Zugriff. In der Praxis zeigt sich jedoch immer wieder, dass schwache, kompromittierte oder mehrfach genutzte Passwörter eine der grössten Sicherheitslücken darstellen. Dieses Skript gibt einen Überblick über die mathematischen Grundlagen der Passwortstärke und fasst die aktuellen Empfehlungen zur Absicherung digitaler Identitäten zusammen \cite{nist2024}.

\subsection{Grundlegende Begriffe und Angriffsarten}
In der Authentisierung und Passwortsicherheit spielen verschiedene Konzepte und Bedrohungen eine entscheidende Rolle:

Neben dem rein mathematischen \gls{bruteforce} nutzen Angreifer in der Praxis sehr häufig ein \gls{dictionaryattack}, bei dem gezielt Listen mit geläufigen Wörtern, Namen und bekannten Daten durchprobiert werden. Um diesen Methoden zu entgehen, ist eine hohe \gls{entropie} entscheidend. Zusätzlich schützt moderne \gls{mfa} (z.\,B. eine Bestätigung direkt auf dem Smartphone) effektiv vor unbefugtem Zugriff.

\subsection{Mathematische Grundlage: Die Passwort-Entropie}
\label{subsec:entropie}
Die Stärke eines \gls{passwort}s gegenüber Offline-Angriffen lässt sich mathematisch über die Shannon-\gls{entropie} beschreiben. Sie gibt an, wie viele Versuche ein Angreifer im schlechtesten Fall benötigt, um den gesamten Suchraum zu durchsuchen.

Die Formel zur Berechnung der Entropie $E$ (in Bit) lautet:
\[
E = L \cdot \log_2(R)
\]

Dabei stehen die Variablen für folgende Werte:
\begin{itemize}
    \item $L$: Die Länge des \gls{passwort}s (Anzahl der verwendeten Zeichen).
    \item $R$: Die Grösse des verfügbaren Zeichensatzes (z.\,B. Kleinbuchstaben, Grossbuchstaben, Zahlen, Sonderzeichen).
\end{itemize}
\newpage
\subsubsection{Beispielrechnung}
\begin{itemize}
    \item \textbf{8 Zeichen (nur Kleinbuchstaben):} Bei $L = 8$ und einem Zeichensatz von $R = 26$ ergibt sich eine Entropie von:
    \[
    E_8 = 8 \cdot \log_2(26) \approx 37{,}6 \text{ Bit}
    \]
    \item \textbf{Verdopplung auf 16 Zeichen (gleicher Zeichensatz):} Bei $L = 16$ und $R = 26$ steigt die Entropie linear mit der Länge auf:
    \[
    E_{16} = 16 \cdot \log_2(26) \approx 75{,}2 \text{ Bit}
    \]
\end{itemize}

In \cref{tab:entropie_vergleich} sind die Entropiewerte für verschiedene Passwortlängen und Zeichensätze zusammenfassend dargestellt.

\begin{table}[htbp]
    \centering
    \begin{tabular}{|l|c|c|r|}
        \hline
        \textbf{Zeichensatz ($R$)} & \textbf{Länge ($L$)} & \textbf{Formel} & \textbf{Entropie ($E$)} \\
        \hline
        Kleinbuchstaben ($R=26$) & 8 & $8 \cdot \log_2(26)$ & $\approx 37{,}6$ Bit \\
        Kleinbuchstaben ($R=26$) & 16 & $16 \cdot \log_2(26)$ & $\approx 75{,}2$ Bit \\
        Alle Zeichen ($R=94$) & 8 & $8 \cdot \log_2(94)$ & $\approx 52{,}4$ Bit \\
        Kleinbuchstaben ($R=26$) & 20 & $20 \cdot \log_2(26)$ & $\approx 94{,}0$ Bit \\
        \hline
    \end{tabular}
    \caption{Vergleich der Passwort-Entropie für verschiedene Längen und Zeichensätze.}
    \label{tab:entropie_vergleich}
\end{table}

\subsubsection{Fazit der Berechnung}
Da die Anzahl der auszuprobierenden Kombinationen im Hintergrund exponentiell mit $2^E$ wächst, schützt das Verlängern des \gls{passwort}s um ein Vielfaches effektiver vor Angriffen als das krampfhafte Einfügen von Sonderzeichen bei einem zu kurzen \gls{passwort}. Zudem verhindert ein langes, zufälliges Passwort, dass ein \gls{dictionaryattack} erfolgreich ist.

\subsection{Aktuelle Richtlinien und Best Practices}
Moderne Richtlinien haben sich von alten Regeln verabschiedet. Heute gelten folgende Hauptempfehlungen:
\begin{enumerate}
    \item \textbf{Länge vor Komplexität:} Bevorzugung von \gls{passphrase}n aus mehreren zufälligen Wörtern. Selbst wenn Angreifer raffinierte Wörterbuchangriffe oder Kombinationen ausprobieren, bricht ein System bei langen Passphrasen (oft auch ohne erzwungene Sonderzeichen) ein, weil die schiere Länge den Suchraum unknackbar macht.
    \item \textbf{Einsatz von \gls{passwortmanager}:} Generierung und Verwaltung unikater Passwörter.
    \item \textbf{Aktivierung von \gls{mfa}:} Einrichtung eines zweiten Faktors (z.\,B. aufs Smartphone).
\end{enumerate}

\newpage
% --- TEIL 2: PHISHING ---
\section{Phishing}
\label{sec:phishing}

\subsection{Einleitung}
In der digitalen Welt bilden nicht nur technische Barrieren wie Passwörter die Abwehrlinie, sondern vor allem der Mensch selbst. Eines der erfolgreichsten und am häufigsten genutzten Einfallstore für Cyberkriminelle ist \gls{phishing}. Dabei wird gezielt die menschliche Psyche ausgenutzt, um an vertrauliche Daten zu gelangen oder Schadsoftware zu verbreiten \cite{bsi2023smtpsmuggling}.

\subsection{Grundlegende Begriffe und Angriffsarten}
\gls{phishing} gehört zu den klassischen \gls{socialengineering}-Angriffen. Statt komplexe technische Systeme direkt zu hacken, manipulieren Angreifer gezielt Personen:
\begin{itemize}
    \item \textbf{Klassisches Phishing:} Täuschend echte E-Mails, SMS oder Nachrichten, die vorgeben, von vertrauenswürdigen Institutionen (wie Banken, Streaming-Diensten oder Behörden) zu stammen.
    \item \textbf{Spear-Phishing:} Ein gezielter Angriff (\gls{spearphishing}), der auf eine spezifische Person oder bestimmte Mitarbeiter zugeschnitten ist, um Vertrauen zu erwecken.
    \item \textbf{Vishing \& Smishing:} Betrugsversuche über Telefonanrufe (\textit{Voice}) oder Kurznachrichten (\textit{SMS}).
    \item \textbf{Quishing:} Das Platzieren von manipulierten QR-Codes (z. B. auf Plakaten oder Parkuhren), die zu gefälschten Webseiten führen.
\end{itemize}

\subsection{Psychologische Faktoren und Wirkungsweise}
Im Gegensatz zu mathematischen Angriffen zielt \gls{phishing} auf menschliche Schwachstellen wie Neugier, Hilfsbereitschaft, Gier oder vor allem \textbf{Angst und Zeitdruck} ab. Typische Maschen beinhalten drohende Konsequenzen wie angebliche Kontosperrungen, dringende Rechnungen oder verpasste Gewinne, damit das Opfer unüberlegt handelt.

\subsection{Aktuelle Richtlinien und Best Practices}
Um sich vor \gls{phishing}-Angriffen zu schützen, gelten im Rahmen der IT-Sicherheit folgende zentrale Empfehlungen:

\begin{enumerate}
    \item \textbf{Kritische Hürden bewahren:} Niemals unüberlegt auf Links klicken oder Anhänge von unbekannten Absendern öffnen.
    \item \textbf{Absender und URLs genau prüfen:} Auf Unstimmigkeiten in E-Mail-Adressen (z. B. Buchstabendreher in Domains) achten.
    \item \textbf{Multi-Faktor-Authentisierung (MFA) nutzen:} Selbst wenn Anmeldedaten abgegriffen werden, verhindert ein zusätzlicher Faktor (wie eine Bestätigung auf dem Smartphone) unbefugten Zugriff.
    \item \textbf{Keine sensiblen Daten preisgeben:} Seriöse Institutionen fragen niemals per E-Mail, SMS oder Telefon nach Passwörtern oder PINs.
\end{enumerate}

\newpage
% --- TEIL 3: ÜBUNGSAUFGABEN ---
\section{Übungsaufgaben}
\label{sec:aufgaben}

Dieser Abschnitt enthält Vertiefungs- und Übungsaufgaben zu den Themen aus \cref{sec:passwortsicherheit} und \cref{sec:phishing}.

\subsection{Aufgabenblatt 1: Passwortsicherheit}
\label{subsec:aufgaben_passwort}

\subsubsection{Teil 1: Begriffsdefinitionen \& Grundlagen}

\begin{taskbox}[{Aufgabe 1.1: Begriffsdefinitionen}]
Erkläre kurz in deinen eigenen Worten den Unterschied zwischen einem Brute-Force-Angriff und einem Dictionary Attack (Wörterbuchangriff).
\end{taskbox}

\begin{taskbox}[{Aufgabe 1.2: Multi-Faktor-Authentisierung (MFA)}]
Was bedeutet die Abkürzung \gls{mfa} und wie schützt sie einen Account, selbst wenn ein Angreifer das Passwort herausgefunden hat?
\end{taskbox}

\subsubsection{Teil 2: Mathematische Berechnungen (Entropie)}

Verwende die Formel zur Berechnung der Entropie aus \cref{subsec:entropie} (siehe auch \cref{tab:entropie_vergleich}):
\[
E = L \cdot \log_2(R)
\]

\begin{taskbox}[{Aufgabe 2.1: Einfluss des Zeichensatzes ($R$)}]
Gegeben ist ein Passwort der Länge $L = 10$.
\begin{enumerate}[label=\alph*)]
    \item Berechne die Entropie $E$, wenn nur Kleinbuchstaben verwendet werden ($R = 26$).
    \item Berechne die Entropie $E$, wenn Klein-, Grossbuchstaben, Zahlen und Sonderzeichen verwendet werden ($R = 94$).
\end{enumerate}
\end{taskbox}

\begin{taskbox}[{Aufgabe 2.2: Länge vs. Komplexität}]
Betrachte die folgenden beiden Passwörter:
\begin{itemize}
    \item \textbf{Szenario A:} Ein komplexes Passwort mit 8 Zeichen ($L = 8$), das Kleinbuchstaben, Grossbuchstaben, Zahlen und Sonderzeichen nutzt ($R = 94$).
    \item \textbf{Szenario B:} Eine einfache Passphrase aus 20 Zeichen ($L = 20$), die nur aus Kleinbuchstaben besteht ($R = 26$).
\end{itemize}
\textbf{Aufgabe:} Berechne für beide Szenarien die Entropie $E$. Welches Passwort bietet eine höhere Entropie und warum entspricht dies der Regel „Länge vor Komplexität“?
\end{taskbox}

\subsubsection{Teil 3: Anwendung \& Best Practices}

\begin{taskbox}[{Aufgabe 3.1: Richtlinien beurteilen}]
Welche der folgenden Aussagen entsprechen den aktuellen Best Practices (z.\,B. nach \gls{nist})? Begründe deine Antwort kurz.
\begin{enumerate}[label=\Alph*)]
    \item „Ein sicheres Passwort muss alle 30 Tage geändert werden.“
    \item „Passphrasen aus mehreren zufälligen Wörtern sind in der Regel sicherer und merkbarer als kurze, komplexe Zeichenketten.“
    \item „Verwende das gleiche Passwort für mehrere Dienste, solange es sehr lang ist.“
    \item „Der Einsatz eines Passwort-Managers wird empfohlen, um für jeden Dienst ein unikat zufälliges Passwort zu generieren.“
\end{enumerate}
\end{taskbox}

\newpage
\subsection{Aufgabenblatt 2: Phishing, Social Engineering \& Best Practices}
\label{subsec:aufgaben_phishing}

\subsubsection{Teil 1: Verständnis \& Begriffe}

\begin{taskbox}[{Aufgabe 1.1: Begriffsdefinitionen \& Abgrenzung}]
Erkläre kurz in deinen eigenen Worten den Unterschied zwischen klassischem Phishing und gezieltem Spear-Phishing. Gehe dabei darauf ein, wer jeweils das Ziel des Angriffs ist.
\end{taskbox}

\begin{taskbox}[{Aufgabe 1.2: Gefahrenkanäle (Vishing, Smishing \& Quishing)}]
Cyberkriminelle nutzen nicht nur E-Mails. Ordne den folgenden Angriffsformen die passende Methode zu und erkläre kurz, worauf man achten muss:
\begin{enumerate}[label=\alph*)]
    \item \textbf{Vishing}
    \item \textbf{Smishing}
    \item \textbf{Quishing}
\end{enumerate}
\end{taskbox}

\begin{taskbox}[{Aufgabe 1.3: Multi-Faktor-Authentisierung (MFA)}]
Warum ist \gls{mfa} eines der effektivsten Schutzmittel gegen Phishing? Wie schützt \gls{mfa} einen Account, selbst wenn ein Opfer seine Anmeldedaten auf einer gefälschten Phishing-Website eingegeben hat?
\end{taskbox}

\newpage
\subsubsection{Teil 2: Psychologische Faktoren \& Szenario-Analyse}

\begin{taskbox}[{Aufgabe 2.1: Psychologische Tricks}]
Social Engineering nutzt primär menschliche Schwachstellen statt technischer Systemlücken. Nenne drei psychologische Faktoren (z.\,B. Angst), die Angreifer gezielt ausnutzen, und beschreibe kurz, wie diese einen Menschen zu unüberlegtem Handeln verleiten.
\end{taskbox}

In \cref{fig:phishing_email} ist eine typische verdächtige Nachricht für die nachfolgende Szenario-Analyse dargestellt.

\begin{figure}[H]
\centering
\begin{emailbox}[Phishing-E-Mail-Verdachtsfall]
\textbf{Von:} \texttt{support-service@bank-security-check.ch} \\
\textbf{An:} \texttt{max.muster@kanti.ch} \\
\textbf{Betreff:} \textbf{DRINGEND: Ihr Konto wird in 2 Stunden gesperrt!} \\[0.5em]
Sehr geehrter Kunde, \\[0.5em]
aufgrund verdächtiger Aktivitäten müssen Sie Ihre Daten sofort verifizieren. Wenn Sie nicht innerhalb von 120 Minuten handeln, wird Ihr Zugang dauerhaft gesperrt. \\[0.5em]
Klicken Sie hier: \url{http://www.kanton-bank-login-sicherheit.com/verify}
\end{emailbox}
\caption{Verdächtiges E-Mail-Fallbeispiel für die Szenario-Analyse.}
\label{fig:phishing_email}
\end{figure}

\begin{taskbox}[{Aufgabe 2.2: Fallbeispiel – Szenario-Analyse}]
Lies die in \cref{fig:phishing_email} dargestellte Nachricht aufmerksam durch.
\begin{enumerate}[label=\alph*)]
    \item Identifiziere mindestens drei Merkmale in dieser Nachricht, die eindeutig auf einen Phishing-Angriff hinweisen.
    \item Welche Form von Phishing und welcher psychologische Mechanismus liegen hier vor?
\end{enumerate}
\end{taskbox}

\subsubsection{Teil 3: Anwendung \& Best Practices (Richtlinien beurteilen)}

\begin{taskbox}[{Aufgabe 3.1: Richtlinien beurteilen}]
Beurteile die folgenden Aussagen. Ist die Aussage \textbf{Richtig} oder \textbf{Falsch}? Begründe deine Antwort jeweils kurz in 1–2 Sätzen.

\textbf{Aussage A:} „Seriöse Banken und Institutionen fordern Kunden gelegentlich per E-Mail auf, ihr Passwort einzugeben, um die Identität bei Sicherheitsupdates zu bestätigen.“

\textbf{Aussage B:} „Beim Aufruf von Links sollte stets die exakte URL (Internetadresse) geprüft werden, da Angreifer häufig Buchstabendreher oder gefälschte Subdomains (Typosquatting) nutzen.“

\textbf{Aussage C:} „Ein Phishing-Angriff kann nur dann erfolgreich sein, wenn auf einem Computer keine Antiviren-Software installiert ist.“

\textbf{Aussage D:} „Das Scannen von QR-Codes auf Plakaten oder Parkuhren ist völlig risikofrei, da QR-Codes keine Viren oder Links enthalten können.“
\end{taskbox}

\newpage
% --- LITERATURVERZEICHNIS & AKRONYME (AUF EINER GEMEINSAMEN SEITE) ---
\printbibliography[heading=bibintoc, title={Literaturverzeichnis}]

\vspace{1.5cm}
\printglossary[type=\acronymtype, title=Akronyme]

\newpage
% --- GLOSSAR (AUF EINER SEPARATEN SEITE) ---
\printglossary[title=Glossar]

\end{document}
